%! Parametric Functions Modeling Planar Motion — Unit 4, Lesson 4.2 — Student Packet Cover Sheet
\documentclass[10pt]{article}
\usepackage{precalculus-article}
\usepackage{precalculus-boxes}
\usepackage{ltablex}
\keepXColumns

\begin{document}

% Full-bleed navy banner
\begin{tikzpicture}[remember picture, overlay]
  \fill[navy] (current page.north west)
    rectangle ([yshift=-0.9in]current page.north east);
\end{tikzpicture}
\vspace{-0.8in}

\begin{center}
  {\color{white}\textbf{\LARGE Precalculus}}\\[4pt]
  {\color{skymid}\large\textbf{Unit 4: Functions Involving Parameters, Vectors, and Matrices}}\\[6pt]
  {\color{white}\textbf{\large Lesson 4.2 \quad Parametric Functions Modeling Planar Motion}}
\end{center}

\vspace{0.22in}
\namedateperiod
\vspace{0.18in}

\begin{learningtargetbox}
\small
\begin{enumerate}[leftmargin=1.5em, itemsep=5pt]
  \item I can identify the horizontal extrema of a parametric motion
    curve by finding the maximum and minimum values of $x(t)$, and the
    vertical extrema by finding the maximum and minimum values of $y(t)$.
  \item I can determine the $x$-intercepts and $y$-intercepts of a
    parametric curve from the zeros of $y(t)$ and $x(t)$, respectively.
  \item I can describe a particle's range of motion by stating the
    maximum and minimum values of each component function over its
    domain.
\end{enumerate}
\end{learningtargetbox}

\vspace{0.14in}

\begin{tocbox}
\small
\renewcommand{\arraystretch}{1.5}
\begin{tabularx}{\linewidth}{@{} c l X r @{}}
\rowcolor{sky}
\textbf{\#} & \textbf{Component} & \textbf{Description} & \textbf{Score} \\
\midrule
1 & Warm-Up        & Spiral Review: Quadratic Extrema \& Component Evaluation & \blank{1.2cm} \\
2 & Group Activity & Remote-Control Car: Extrema and Intercepts & \blank{1.2cm} \\
3 & Exit Ticket    & Planar Motion Check                                        & \blank{1.2cm} \\
4 & Homework       & Parametric Planar Motion Practice                          & \blank{1.2cm} \\
\midrule
  & \multicolumn{2}{r}{\textbf{Total}} & \blank{1.2cm} \\
\end{tabularx}
\end{tocbox}

\end{document}
